% ----------------------------------------------------------------
\documentclass{article}
\usepackage[svgnames]{xcolor}
\typeout{about to load package}
\usepackage{atableau2}
\typeout{package loaded}
\usepackage{enumitem}
\setlist[description]{font=\sffamily\bfseries\color{NavyBlue},labelwidth=\textwidth}
\usepackage[a4paper,margin=18mm]{geometry}
\synctex=1
\parindent=0pt
\parskip=2mm

\newcommand\heading[1]{\bigskip\noindent\textsf{\large #1}}

\newcommand\Section[1]{\subsection{\textcolor{NavyBlue}{#1}}}

\ExplSyntaxOn
\NewDocumentCommand\option{ omo }{%
  \IfNoValueTF{#1}{\textsf{#2}}
  {%
    \textbf{\textcolor{Crimson}{#2}}
    \str_if_empty:nF {#1}{~(default:~\textcolor{DarkRed}{#1})}%
    \IfNoValueF{#3}{\hfill\textcolor{Gray}{[accepts:~#3]}}%
  }%
  \index{#2}%
}
\ExplSyntaxOff

% ----------------------------------------------------------------
\usepackage{imakeidx}
\indexsetup{level=\section*, toclevel=section, noclearpage}
\makeindex[intoc,columns=3]

% ----------------------------------------------------------------
\usepackage[skins,listings]{tcolorbox}
\usetikzlibrary{backgrounds, shapes.geometric} % for tikz={framed,...}

\DeclareTotalTCBox\keyword{ O{} v }{
  fontupper=\sffamily,
  nobeforeafter,
  skin=tile,
  verbatim,
  on line,
  tcbox raise base,
  top=0pt,bottom=0pt,left=0mm,right=0mm,
  colback=OldLace,
  colupper=RoyalBlue!70,
  #1}
{#2}

\lstdefinestyle{tikz}{style=tcblatex,
  classoffset=0,
  texcsstyle=*\color{DarkGoldenrod},%
  deletetexcs={begin, end},
  moretexcs={,%
    Tableau,
    Tabloid,
    ShiftedTableau,
    Diagram,
    ShiftedDiagram,
    Abacus,
  },%
  classoffset=1,
  keywordstyle=\color{NavyBlue},%
  morekeywords={wordle,GrilleSutom},
  classoffset=2,
  keywordstyle=\color{Crimson},%
  morekeywords={
    Australian,
    English,
    French,
    Ukraine,
    australian,
    english,
    finite,
    font,
    french,
    infinite,
    shape,
    scale,
    yscale,
    xscale,
    shifted,
    tabloid,
    ukraine,
  },
  classoffset=3,
  keywordstyle=\color{DarkOrange},%
  morekeywords={
    tikzset,
    atabset
  }
}

\DeclareTCBListing{example}{ !O{} }{%
  skin=bicolor,
  colframe=RoyalBlue!50,
  colbacklower=OldLace,
  colback=LightSkyBlue!20,
  lefthand width=0.35\textwidth,
  listing style=tikz,
  sidebyside,
  sidebyside align=center,
  sidebyside gap=4mm,
  text and listing,
  text outside listing,
  boxsep = 0pt,
  #1
}


% ----------------------------------------------------------------
% manual title
\makeatletter
\author{Andrew Mathas}
\usepackage{tikz}
\usetikzlibrary{shadows.blur}
\tikzset{
  shadowed/.style={
    blur shadow={shadow blur steps=5},
    bottom color=LightSkyBlue!30,
    draw=NavyBlue!70,
    shade,
    font=\normalfont\Huge\bfseries\scshape,
    rounded corners=8pt,
    top color=RoyalBlue,
  },
  boxes/.style={draw=RoyalBlue,
    fill=AliceBlue,
    font=\sffamily\small,
    inner sep=5pt,
    rectangle,
    rounded corners=8pt,
    text=MidnightBlue,
  }
}
\newcommand\ATableau{%
  \begin{tikzpicture}[remember picture,overlay]
      \node[yshift=-3cm] at (current page.north west)
        {\begin{tikzpicture}[remember picture, overlay]
          \draw[shadowed](30mm,0) rectangle node[white]{aTableau} (\paperwidth-30mm,16mm);
          \node[anchor=west,boxes] at (4cm,0cm) {\@author};
          \node[anchor=east,boxes] at (\paperwidth-4cm,0) {Version \atableau@version};
         \end{tikzpicture}
        };
   \end{tikzpicture}
   \vspace*{20mm}
}

\def\@oddfoot{\textsc{Wordle} --- \atableau@version~(released~\atableau@release)\hfill\thepage}

\usepackage[colorlinks=true,linkcolor=blue,urlcolor=MediumBlue]{hyperref}
\hypersetup{
  pdfcreator={ Generated by pdfLaTeX },
  pdfinfo={
    Author  ={ Andrew Mathas },
    Keywords={ algebraic combinatorics, tableau, abacus, combinatorics,  },
    License ={ LaTeX Project Public License v1.3c or later },
    Subject ={ LaTeXing Wordle puzzles },
    Title   ={ ATableau - \atableau@version }
  },
}
\makeatother


\begin{document}

  \ATableau

  \tableofcontents

  This package provides commands for drawing Young diagrams and
  tableaux, which are

\section{Tableaux}

  \begin{example}
    \Tableau{123,45}
  \end{example}

  \begin{example}
    \begin{tikzpicture}
      \node at (1,1) {$\bullet$};
      \node at (2,2.5) {$\bullet$};
      \Tableau(1,1) {134,5,.7}
      \Tableau(2,2.5) {{10}{11},19,8}
    \end{tikzpicture}
  \end{example}

  \begin{example}
    \Tableau[starstyle={fill=blue}]{..*., .*.*., .*.}
    \Tableau[dotstyle={fill=blue!20}]{..., ..., ..}
    \Tableau[dot=x]{..., ..., ..}
  \end{example}
  \begin{example}
    \tikzset{R/.style={fill=blue!20,draw=red,thick}}
    \Tableau{12[R]3,4[R]5}
    \Tableau[starstyle={fill=blue!50,text=white}]{12*3,4*5}
  \end{example}

  \begin{example}
    \tikzset{R/.style={fill=red!20,draw=red}}
    \Tableau[ukraine]{12[R]3,45}
  \end{example}

  \begin{example}
    \tikzset{R/.style={fill=blue!20,draw=NavyBlue}}
    \Tableau[Australian]{12[R]3,45}
  \end{example}

  \begin{example}
    \tikzset{R/.style={fill=red!20,circle, radius=0.2}}
    \Tableau[french]{12[R]3,45}
  \end{example}

  \begin{example}
    \tikzset{R/.style={fill=red!20,circle}}
    \ShiftedTableau{12[R]3,45}
  \end{example}

  \begin{example}
    \tikzset{R/.style={fill=red!20,circle}}
    \ShiftedTableau[french]{12[R]3,45}
  \end{example}

  \begin{example}
    \tikzset{R/.style={fill=red!20,circle}}
    \SkewTableau{2,1}{12[R]3,45,67}
    \SkewTableau[french]{2,1}{12[R]3,45,67}
  \end{example}

\end{document}
  \heading{Abacuses}

  Use the \LaTeXCode{\EAbacus} command to a typeset an abacus

  \begin{example}
     Usage:  \textbackslash EAbacus{e}{rows}{partition}
     \EAbacus{3}{4}{5,4,1,1}   % \EAbacus{3}{4}{5,4,1,1}
  \end{example}

  \heading{Braid diagrams}

  Use the \LaTeXCode{\braid} command to typeset briad diagrams:

  \begin{example}
       \braid{1, 4, 2, 3}     % \braid{1, 4, 2, 3}
       \braid*{3, 1, 4, 2}    % \braid*{3, 1, 4, 2}
       \braid!{2, 3, 4, 1}    % \braid!{2, 3, 4, 1}
       \braid*!{4, 3, 2, 1}   % \braid*!{4, 3, 2, 1}
  \end{example}

  \[
         \braid[scale=2]*!{4,3,2,1} =
  \]

                \braid*{3,2,4,1}\braid!{2,3,1,4}

  \heading{Taboids}

  Use the \LaTeXCode{\Tableau} command to typeset a tabloids:

  \begin{example}
    \Tabloid{{1,3,7,9,11},{2,4,9,10},{6},{8}}   % \Tabloid{{1,3,7,9,11},{2,4,9,10},{6},{8}}
  \end{example}

  \heading{Tableau}

  Use the \LaTeXCode{\Tableau} command to a typeset tableau:

  \begin{example}
    \Tableau{{1,3,7,9,11},{2,4,9,10},{6},{8}}   % \Tableau{{1,3,7,9,11},{2,4,9,10},{6},{8}}
  \end{example}

  \begin{example}
      \Tableau{ {,\bullet,,,,},{\bullet,B,\bullet,},{,\bullet,,},{,},{}}
      %\Tableau{ {,\bullet,,,,},{\bullet,B,\bullet,},{,\bullet,,},{,},{}}
  \end{example}

  \heading{Young diagrams}

  Use the \LaTeXCode{\Diagram} command to typeset Young diagrams:

  \begin{example}
       \Diagram{5,4,1,1}     % \Diagram{5,4,1,1} %
  \end{example}


\end{document}
